\section{Euler--Kronecker constant of the character of dodecaphase type}

The factorization of Dedekind

\begin{equation}
\zeta_{\mathbb Q(\sqrt3)}(s)=\zeta(s)L(s,\chi_{12})
\label{eq:dedekind-sqrt3}
\end{equation}

gives, upon comparing the constant terms in \(s=1\),

\begin{equation}
\gamma_{\mathbb Q(\sqrt3)}
=\gamma+\frac{L'(1,\chi_{12})}{L(1,\chi_{12})}
=1.0539656082484825800\ldots.
\label{eq:euler-kronecker}
\end{equation}

This constant represents the normalized finite term at the pole of the
zeta function of the real dodecaphase field. It is not a variant of
Euler's \(\gamma\), but the arithmetic correction induced by
\(\chi_{12}\).

In fact, writing \(u=s-1\),
\[
\zeta(s)=u^{-1}+\gamma+O(u),
\qquad
L(s,\chi_{12})=L_1+L_1'u+O(u^2),
\]
we obtain
\begin{equation}
\zeta_{\mathbb Q(\sqrt3)}(s)
=\frac{L_1}{u}+\bigl(\gamma L_1+L_1'\bigr)+O(u).
\label{eq:ek-laurent-derivation}
\end{equation}
The Euler--Kronecker constant is the quotient of the finite term by
the residue; dividing \cref{eq:ek-laurent-derivation} by \(L_1\) proves
\cref{eq:euler-kronecker}. This derivation simultaneously explains the
residue
\[
\operatorname*{Res}_{s=1}\zeta_{\mathbb Q(\sqrt3)}(s)
=\frac{\log(2+\sqrt3)}{\sqrt3}
\]
and the finite correction. The numerical value is certified by evaluating
\(L(1,\chi_{12})\) and \(L'(1,\chi_{12})\) with the same residue table and a
tail bound; it is not imported from a table of quadratic fields.

\section{Constantes aritmeticas asociadas to the numeros primos}

Once the irreducible arithmetic cycles have been selected, two
distinct families are obtained:

\begin{align}
B_1&=\lim_{x\to\infty}\left(
\sum_{p\le x}\frac1p-\log\log x
\right),
\label{eq:meissel-mertens}\\
C_2&=\prod_{p>2}\frac{p(p-2)}{(p-1)^2}.
\label{eq:twin-prime-constant}
\end{align}

\(B_1\) regularizes a local sum; \(C_2\) is a multiplicative density of twin-prime pairs. Both require tail control. They are not deduced from one another by sharing the primes.

